\section{La divergencia profunda entre 2C y 2B}

\subsection{El estancamiento de percepción en el lado 2C}

Yao Shunyu (姚顺雨) observó un fenómeno interesante: cuando la gente piensa en un AI Super App, lo que ahora viene a la mente son dos productos: ChatGPT y Claude Code, que representan respectivamente los ejemplos típicos de 2C y 2B.

\begin{warningbox}{La paradoja del lado 2C}
Para la mayoría de las personas, la experiencia de usar ChatGPT hoy ha cambiado poco respecto a hace un año. Aunque capacidades matemáticas avanzadas, como el álgebra abstracta o la teoría de Galois, mejoraron de manera notable, el usuario común no lo percibe. En China, la mayoría todavía considera la IA como un «motor de búsqueda mejorado», sin saber cómo despertar la verdadera inteligencia del modelo.
\end{warningbox}

\subsection{El valor lineal del lado 2B}

\begin{importantbox}{La lógica de valor del lado 2B}
En el lado 2B, mayor inteligencia = mayor productividad = mayores ingresos. Estos objetivos están altamente aligned entre sí. Por ejemplo, un modelo del nivel de OPUS 4.5 podría resolver correctamente de manera directa 8--9 de cada 10 tareas, mientras que un modelo peor solo acierta 5--6. El usuario no sabe cuáles se resolvieron mal y necesita invertir mucho esfuerzo en supervisarlas. Por eso, los usuarios empresariales están dispuestos a pagar una prima (por ejemplo, 200 dólares al mes) por usar el mejor modelo.
\end{importantbox}

La dificultad de hacer IA para 2C radica en que indicadores de producto como el DAU a menudo no se correlacionan con la inteligencia del modelo, e incluso pueden ser inversamente proporcionales. En cambio, en el lado 2B, mientras mejor es el modelo, mayores son los ingresos: todo está aligned.

\subsection{El verdadero cuello de botella del lado 2C: el context, no el modelo}

\begin{knowledgebox}{Un ejemplo revelador}
«¿Qué debería comer hoy?»: esta pregunta, ya sea que se le haga a ChatGPT este año, el año pasado o el próximo, no obtiene una buena respuesta. Porque el cuello de botella no está en un modelo más grande o un RL más fuerte, sino en \textbf{el context adicional}: si hace frío o no, el rango de actividades, las preferencias de la familia, etc. Si se pudiera aprovechar bien un insumo adicional como el historial de chat de WeChat, eso aportaría más valor al usuario.
\end{knowledgebox}

\subsection{Resumen de la sección}

El cuello de botella del lado 2C está en el context y el environment, no en el modelo en sí; la cadena de valor del lado 2B es altamente aligned y representa una dirección de crecimiento más segura.

